\section{Representation, Learnability, and the Virtue of Complexity}
\label{sec:representation_learnability}
\citet{kellymalamudzhou2024} study an environment in which expanding the model
also captures an increasing share of the underlying economic signal. In their
setting, the improvement in approximation can dominate the additional
estimation cost. This is the virtue of complexity.

Our question is different. We hold the investor's information set fixed and
vary only the richness of the portfolio representation. We take the view that
a portfolio manager should make available all information she considers
economically useful and can observe in real time. Regularization, rather than
ex ante feature exclusion, should then determine how much of the resulting
portfolio opportunity set can actually be learned from the available data.
We therefore study the virtue of Complexity, as the virtue of richer representation, separately from the
addition of new signals.


Our point in this section is that a richer portfolio representation and a
better learned portfolio are two different objects.

Let $\SR^\star$ denote the maximum population Sharpe ratio over the space of
$L^2$ portfolio policies,
\[
\SR^\star := \sup_{W\in L^2}\SR(W),
\]
and let
\[
\SR_K^\star := \sup_{W\in\mathcal H_K}\SR(W)
\]
denote the maximum population Sharpe ratio available within representation
$K$, where $\mathcal H_K\subseteq L^2$ is the corresponding policy space.
For example, a linear kernel restricts portfolio weights to be linear
functions of characteristics, whereas a Matérn kernel allows a richer
nonlinear class.

Finally,
\[
\SR(\widehat W_{K,\lambda,T})
\]
denotes the population Sharpe ratio of the portfolio policy learned within
representation $K$ from a market history of length $T$.

The gap between the unrestricted population opportunity set and the portfolio
actually learned from the data can be decomposed as
\begin{equation}
\SR^\star-\SR(\widehat W_{K,\lambda,T})
=
\underbrace{\SR^\star-\SR_K^\star}_{\text{representation gap}}
+
\underbrace{
\SR_K^\star-\SR(\widehat W_{K,\lambda,T})
}_{\text{learning gap}}.
\label{eq:representation_estimation_decomposition}
\end{equation}

The representation gap measures the investment opportunities that cannot be
expressed by the chosen representation, even if the investment environment
were known perfectly. The learning gap measures how much of the opportunity
set available within that representation cannot be recovered from a finite
history of returns.

Now consider two nested representations,
\[
\mathcal H_{K_1}\subseteq\mathcal H_{K_2}.
\]
At the population level,
\[
\SR_{K_2}^\star\geq\SR_{K_1}^\star.
\]
A richer representation can therefore only expand the population investment
opportunity set. In this precise sense, complexity is a virtue for
representation.

But this population ordering need not survive estimation from finite data. It
is perfectly possible that
\[
\SR_{K_2}^\star>\SR_{K_1}^\star
\]
while
\[
\SR(\widehat W_{K_2,\lambda_2,T})
<
\SR(\widehat W_{K_1,\lambda_1,T}).
\]
The larger representation contains better population opportunities, but some
of those opportunities may lie in managed-portfolio directions that are too
weakly identified by the available return history.

A richer representation therefore expands what the portfolio manager could
exploit. Whether those additional opportunities improve the learned portfolio
depends on whether they can actually be learned from the available data.